\beginchapter Chapter 18. Fine Points of\\Mathematics\\Typing

We have discussed most of the facilities needed to construct math
formulas, but there are several more things a good mathematical typist
will want to watch for. After you have typed a dozen or so formulas using
the basic ideas of Chapters 16 and~17, you will find that it's easy to
visualize the final appearance of a mathematical expression as you type
it.  And once you have gotten to that level, there's only a little bit
more to learn before you are producing formulas as beautiful as any the
world has ever seen; tastefully applied touches of \TeX nique will add a
professional polish that works wonders for the appearance and readability
of the books and papers that you type. This chapter talks about such
tricks, and it also fills in a few gaps by mentioning some aspects of math
that didn't fit comfortably into~Chapters~16~and~17.

\subsection Punctuation.  When a formula is followed by a ^{period}, ^{comma},
^{semicolon}, ^{colon}, ^{question mark}, ^{exclamation point}, etc., put the
^{punctuation} {\sl after\/} the |$|, when the formula is in the text; but
put the punctuation {\sl before\/} the |$$| when the formula is displayed.
For example,
\begintt
If $x<0$, we have shown that $$y=f(x).$$
\endtt
\TeX's spacing rules within paragraphs work best when the
^{punctuation marks} are not considered to be part of the formulas.

Similarly, don't ever type anything like
\begintt
for $x = a, b$, or $c$.
\endtt
It should be
\begintt
for $x = a$, $b$, or $c$.
\endtt
(Better yet, use a ^{tie}: `|or~$c$|'.) \ The reason is that \TeX\ will
typeset expression `|$x|~|=|~|a,|~|b$|' as a single formula, so it will
put a ``^{thin space}'' between the comma and the $b$. This space will
not be the same as the space that \TeX\ puts after the
comma {\sl after\/} the $b$, since spaces between words are always bigger than
thin spaces. Such unequal spacing looks bad, but when you type things right
the spacing will look good.

Another reason for not typing `|$x| |=| |a,| |b$|' is that it inhibits the
possibilities for breaking lines in a paragraph: \TeX\ will never break at
the space between the comma and the |b| because breaks after commas in
formulas are usually wrong. For example, in the equation
`|$x|~|=|~|f(a,|~|b)$|' we certainly don't want to put `$x=f(a,$' on one
line and `$b)$' on the next.

Thus, when typing formulas in the text of a paragraph, keep the math properly
segregated: Don't take operators like $-$ and $=$ outside of the |$|'s,
and keep commas inside the formula if they are truly part of the formula.
But if a comma or period or other punctuation mark belongs linguistically
to the sentence rather than to the formula, leave it outside the |$|'s.

\exercise Type this: $R(n,t)=O(t^{n/2})$, as $t\to0^+$.
\answer |$R(n,t)=O(t^{n/2})$, as $t\to0^+$.| \
(N.B.: `|O(|', not `|0('|.)

\danger Some mathematical styles insert a bit of extra space around
formulas to separate them from the text. For example, when copy is
being produced on an ordinary typewriter that doesn't have italic
letters, the best technical typists have traditionally put an extra
blank space before and after each formula, because this provides a
useful visual distinction. You might find it helpful to think of
each |$| as a symbol that has the potential of adding a little space
to the printed output; then the rule about excluding sentence
punctuation from formulas may be easier to remember.

\ddanger \TeX\ does, in fact, insert additional ^{space} before and after each
formula; the amount of such space is called ^|\mathsurround|, which is
a \<dimen>-valued parameter. For example, if you set |\mathsurround=1pt|,
each formula will effectively be 2~points wider ($1\pt$ at each side):
$$\baselineskip13pt\halign{\indent\mathsurround=#pt
  For $x=a$, $b$, or $c$.\hfil&\quad(#)\hfil\cr
1&|\mathsurround=1pt|\cr 0&|\mathsurround=0pt|\cr}$$
This extra space will disappear into the left or right margin if the formula
occurs at the beginning or end of a line. The value of\/ |\mathsurround| that
is in force when \TeX\ reads the closing |$| of a formula is used at both
left and right of that formula.  Plain \TeX\ takes |\mathsurround=0pt|, so
you won't see any extra space unless you are using some other format, or
unless you change |\mathsurround| yourself.

\subsection Non-italic letters in formulas. The names of algebraic variables
are usually italic or Greek letters, but common mathematical functions
like `log' are always set in ^{roman type}. The best way to deal with such
constructions is to make use of the following 32~control sequences (all of
which are defined in plain \TeX\ format, see Appendix~B):
\begintt
\arccos  \cos   \csc   \exp   \ker     \limsup  \min  \sinh
\arcsin  \cosh  \deg   \gcd   \lg      \ln      \Pr   \sup
\arctan  \cot   \det   \hom   \lim     \log     \sec  \tan
\arg     \coth  \dim   \inf   \liminf  \max     \sin  \tanh
\endtt
^^|\arccos| ^^|\cos| ^^|\csc| ^^|\exp| ^^|\ker| ^^|\limsup| ^^|\min|
^^|\sinh| ^^|\arcsin| ^^|\cosh| ^^|\deg| ^^|\gcd| ^^|\lg| ^^|\ln| ^^|\Pr|
^^|\sup| ^^|\arctan| ^^|\cot| ^^|\det| ^^|\hom| ^^|\lim| ^^|\log| ^^|\sec|
^^|\tan| ^^|\arg| ^^|\coth| ^^|\dim| ^^|\inf| ^^|\liminf| ^^|\max|
^^|\sin| ^^|\tanh|
These control sequences lead to roman type with appropriate
spacing:\def\sep{&\hskip-1em}
\beginlongmathdemo
\it Input\sep\it Output\cr
\noalign{\vskip2pt}
|$\sin2\theta=2\sin\theta\cos\theta$|\sep\sin2\theta=2\sin\theta\cos\theta\cr
|$O(n\log n\log\log n)$|\sep O(n\log n\log\log n)\cr
|$\Pr(X>x)=\exp(-x/\mu)$|\sep\Pr(X>x)=\exp(-x/\mu)\cr
|$$\max_{1\le n\le m}\log_2P_n$$|\sep
  \displaystyle{\max_{1\le n\le m}\log_2P_n}\cr
\noalign{\vskip2pt}
|$$\lim_{x\to0}{\sin x\over x}=1$$|\sep\displaystyle{\lim_{x\to0}
  {\sin x\over x}=1}\cr
\endmathdemo
^^|\mu|
The last two formulas, which are displays, show that some of the special
control sequences are treated by \TeX\ as ``large operators'' with limits
just like $\sum$: The subscript on |\max| is not treated like the subscript
on |\log|. Subscripts and superscripts will become limits
when they are attached to |\det|, |\gcd|, |\inf|, |\lim|, |\liminf|,
|\limsup|, |\max|, |\min|, |\Pr|, and |\sup|, in display style.

\exercise Express the following display in plain \TeX\ language, using
`^|\nu|' for `$\nu$':
$$p_1(n)=\lim_{m\to\infty}\sum_{\nu=0}^\infty
  \bigl(1-\cos^{2m}(\nu!^n\pi/n)\bigr).$$
\answer |$$p_1(n)=\lim_{m\to\infty}\sum_{\nu=0}^\infty|\parbreak
|  \bigl(1-\cos^{2m}(\nu!^n\pi/n)\bigr).$$|\par
\smallskip\noindent $\bigl[$Mathematicians may enjoy
interpreting this formula; cf.~G.~H. ^{Hardy},
{\sl Messenger of Mathematics\/ \bf35} (1906), 145--146.$\bigr]$

\danger If you need roman type for some mathematical function or operator
that isn't included in plain \TeX's list of~32, it is easy to define a
new control sequence by mimicking the definitions in Appendix~B\null. Or, if
you need roman type just for a ``one shot'' use, it is even easier to
get what you want by switching to ^|\rm| type, as follows:
\beginlongmathdemo
|$\sqrt{{\rm Var}(X)}$|&\sqrt{{\rm Var}(X)}\cr
|$x_{\rm max}-x_{\rm min}$|&x_{\rm max}-x_{\rm min}\cr
|${\rm LL}(k)\Rightarrow{\rm LR}(k)$|&{\rm LL}(k)\Rightarrow{\rm LR}(k)\cr
|$\exp(x+{\rm constant})$|&\exp(x+{\rm constant})\cr
|$x^3+{\rm lower\ order\ terms}$|&x^3+{\rm lower\ order\ terms}\cr
\endmathdemo
Notice the uses of `|\|\]' ^^{control space} in the last case;
without them, the result would have been `$x^3+{\rm lower order terms}$',
because ordinary blank ^{spaces} are ignored in math mode.

\danger You can also use ^|\hbox| instead of\/ |\rm| to get roman letters
into formulas. For example, four of the last five formulas can be
generated by
\beginlongmathdemo
|$\sqrt{\hbox{Var}(X)}$|&\sqrt{\hbox{Var}(X)}\cr
|$\hbox{LL}(k)\Rightarrow\hbox{LR}(k)$|&\hbox{LL}(k)\Rightarrow\hbox{LR}(k)\cr
|$\exp(x+\hbox{constant})$|&\exp(x+\hbox{constant})\cr
|$x^3+\hbox{lower order terms}$|&x^3+\hbox{lower order terms}\cr
\endmathdemo
In this case `|\|\]' isn't necessary, because the material in an
|\hbox| is processed in horizontal mode, when spaces are significant.
But such uses of\/ |\hbox| have two disadvantages: \ (1)~The contents of the
box will be typeset in the same size, whether or not the box occurs as a
subscript; for example, `|$x_{\hbox{max}}$|' yields `$x_{\hbox{max}}$'.
\ (2)~The font that's used inside |\hbox| will be the ``^{current font},''
so it might not be roman. For example, if you are typesetting the statement
of some theorem that is in slanted type, and if that theorem refers
to `|$\sqrt{\hbox{Var}(X)}$|', you will get the unintended result
`{\sl$\sqrt{\hbox{Var}(X)}$}'. In order to make sure that an |\hbox| uses
roman type, you need to specify |\rm|, e.g., `|$\sqrt{\hbox{\rm Var}(X)}$|';
and then the |\hbox| serves no purpose. We will see later, however, that
|\hbox| can be very useful in displayed formulas.

\ddangerexercise When the displayed formula
`|$$\lim_{n\to\infty}x_n {\rm\ exists} \iff|\break
|\limsup_{n\to\infty}x_n = \liminf_{n\to\infty}x_n.$$|' is typeset with
the standard macros of plain \TeX, you get
$$\lim_{n\to\infty}x_n{\rm\ exists}\iff
  \limsup_{n\to\infty}x_n=\liminf_{n\to\infty}x_n.$$
But some people prefer a different notation: Explain how you could change
the definitions of\/ ^|\limsup| and ^|\liminf| so that the display would be
$$
\def\limsup{\mathop{\overline{\rm lim}}}
\def\liminf{\mathop{\underline{\rm lim}}}
\lim_{n\to\infty}x_n{\rm\ exists}\iff
  \limsup_{n\to\infty}x_n=\liminf_{n\to\infty}x_n.$$
\answer |\def\limsup{\mathop{\overline{\rm lim}}}|\parbreak
|\def\liminf{\mathop{\underline{\rm lim}}}|\par
\smallskip\noindent
[Notice that the limits `$n\to\infty$' appear at different levels, in both
of the displays, because `sup' and the underbar descend below the baseline.
It is possible to unify the limit positions by using ^{phantoms}, as explained
later in this chapter. For example,
\begintt
\def\limsup{\mathop{\vphantom{\underline{}}\overline{\rm lim}}}
\endtt
would give lower limits in the same position as |\liminf|.]

\danger The word `mod' is also generally set in roman type, when it occurs
in formulas; but this word needs more care, because it is used in two
different ways that require two different treatments.
Plain \TeX\ provides two different control sequences,
^|\bmod| and ^|\pmod|, for the two cases: |\bmod| is to be used when
`mod' is a ^{binary operation} (i.e., when it occurs between two quantities,
like a plus sign usually does), and |\pmod| is to be used when
`mod' occurs parenthetically at the end of a formula. For example,
\beginmathdemo
|$\gcd(m,n)=\gcd(n,m\bmod n)$|&\gcd(m,n)=\gcd(n,m\bmod n)\cr
|$x\equiv y+1\pmod{m^2}$|&x\equiv y+1\pmod{m^2}\cr
\endmathdemo
The `|b|' in `|\bmod|' stands for ``binary''; the `|p|' in `|\pmod|' stands
for ``parenthesized.'' Notice that |\pmod| inserts its own parentheses;
the quantity that appears after `mod' in the parentheses should be
enclosed in braces, if it isn't a single symbol.

\dangerexercise What did poor B. L. ^{User} get when he typed
`|$x\equiv0 (\pmod y^n)$|'\thinspace?
\answer $x\equiv0(\pmod y^n)$. He should have typed
`|$x\equiv0\pmod{y^n}$|'.

\dangerexercise Explain how to produce \lower12pt\null\
$\smash{\displaystyle{n\choose k}\equiv{\lfloor n/p\rfloor\choose
\lfloor k/p\rfloor}{n\bmod p\choose k\bmod p}\pmod p.}$
\answer |$${n\choose k}\equiv{\lfloor n/p\rfloor\choose|\parbreak
|  \lfloor k/p\rfloor}{n\bmod p\choose k\bmod p}\pmod p.$$|

\danger The same mechanism that works for roman type in formulas can be used
to get other styles of type as well. For example, ^|\bf| yields ^{boldface}:
\beginmathdemo
|$\bf a+b=\Phi_m$|&\bf a+b=\Phi_m\cr
\endmathdemo
Notice that whole formula didn't become emboldened in this example; the
`$+$' and `$=$' stayed the same. Plain \TeX\ sets things up so
that commands like |\rm| and |\bf| will affect only the uppercase letters
|A|~to~|Z|, the lowercase letters |a|~to~|z|, the digits |0|~to~|9|,
the uppercase Greek letters |\Gamma| to~|\Omega|, and math ^{accents}
like ^|\hat| and ^|\tilde|.  Incidentally, no braces were used in this
example, because |$|'s have the effect of grouping; |\bf| changes the
current font, but the change is local, so it does not affect the font that
was current outside the formula.

\ddanger The bold fonts available in plain \TeX\ are ``bold roman,'' rather
than ``bold italic,'' because the latter are rarely needed. However, \TeX\
could readily be set up to make use of bold math italics, if desired
(see exercise 17.\bmiexno). A more extensive set of math fonts would also
include ^{script}, ^{Fraktur}, and ``^{blackboard bold}'' styles; plain
\TeX\ doesn't have these, but other formats like \AmSTeX\ do. ^^{AMS-TeX}
^^{German black letters}

\danger Besides |\rm| and |\bf|, you can say ^|\cal| in formulas to get
uppercase letters in a ``^{calligraphic}'' style. For example, `|$\cal
A$|' produces `$\cal A$' and `|$\cal Z$|' produces `$\cal Z$'. But beware:
This works only with the letters |A| to |Z|; you'll get weird results if
you apply |\cal| to lowercase or Greek letters.

\danger There's also ^|\mit|, which stands for ``^{math italic}.'' This
affects ^{uppercase Greek}, so that you get
$\mit(\Gamma,\Delta,\Theta,\Lambda,\Xi,\Pi,\Sigma,\Upsilon,\Phi,\Psi,\Omega)$
instead of $(\Gamma,\ldots,\Omega)$.  When~|\mit| is in effect, the
ordinary letters |A| to |Z| and |a| to |z| are not changed; they are set
in italics as usual, because they ordinarily come from the math italic
font. Conversely, uppercase Greek letters and math accents are unaffected
by |\rm|, because they ordinarily come from the roman font. Math accents
should not be used when the |\mit| family has been selected, because the
math italic font contains no accents.

\dangerexercise Type the formula $\bf\bar x^{\rm T}Mx={\rm0}\iff x=0$,
using as few keystrokes as possible.  ^^{boldface numbers in math}
\ (The first `0' is roman, the second is bold. The superscript `T' is roman.)
\answer |$\bf\bar x^{\rm T}Mx={\rm0}\iff x=0$|.  \ (If you typed a space between
|\rm| and~|0|, you wasted a keystroke; but don't feel guilty about it.)

\dangerexercise Figure out how to typeset
`$S\subseteq\mit\Sigma\iff S\in\cal S$'.
\answer |$S\subseteq{\mit\Sigma}\iff S\in{\cal S}$|. In this case the
braces are redundant and could be eliminated; but you shouldn't try to do
{\sl everything\/} with fewest keystrokes, or you'll outsmart yourself
some day.

\danger Plain \TeX\ also allows you to type ^|\it|, ^|\sl|, or ^|\tt|, if
you want text italic, slanted, or typewriter letters to occur in a math
formula. However, these fonts are available only in text size, so you
should not try to use them in subscripts.

\danger If you're paying attention, you probably wonder why both
|\mit| and |\it| are provided; the answer is that |\mit| is ``math italic''
(which is normally best for formulas), and |\it| is ``^{text italic}'' (which
is normally best for running text).
\beginmathdemo
|$This\ is\ math\ italic.$|&This\ is\ math\ italic.\cr
|{\it This is text italic.}|&\hbox{\it This is text italic.}\cr
\endmathdemo
The math italic letters are a little wider, and the spacing is different;
this works better in most formulas, but it fails spectacularly when
you try to type certain italic words like `$different$' using math mode
(`|$different$|'). A wide `$f$' is usually desirable in formulas, but it
is undesirable in text. Therefore wise typists
use |\it| in a math formula that is supposed
to contain an actual italic word. Such cases almost never occur in
classical mathematics, but they are common when ^{computer programs}
are being typeset, since programmers often use multi-letter ``^{identifiers}'':
\beginmathdemo
|$\it last:=first$|&\it last:=first\cr
|$\it x\_coord(point\_2)$|&\it x\_coord(point\_2)\cr
\endmathdemo
The first of these examples shows that \TeX\ recognizes the ^{ligature}
`{\it fi\/}' when text italic occurs in a math formula;
the other example illustrates the use of short ^{underlines} to break
up identifier names. ^^{control-underline}
When the author typeset this manual, he used `|$\it SS$|' to refer to
style~$\SS$, since `|$SS$|' makes the $S$'s too far apart: $SS$.

\dangerexercise What plain \TeX\ commands will produce the following display?
$$\tenmath
{\it available}+\sum_{i=1}^n\max\bigl({\it full}(i),{\it reserved}(i)\bigr)
  ={\it capacity}.$$
\answer |$${\it available}+\sum_{i=1}^n\max\bigl({\it full}(i),|\parbreak
        |{\it reserved}(i)\bigr)={\it capacity}.$$|
\smallskip\noindent [If\/ |\it| had been used throughout
the formula, the subscript~$i$ and superscript~$n$ would have caused error
messages saying `^|\scriptfont| |4| |is| |undefined|',
since plain \TeX\ makes |\it| available only in text size.]

\ddangerexercise How would you go about typesetting the following computer
program, using the macros of plain \TeX?
$$\vbox{\let\par=\endgraf
\obeylines\sfcode`;=3000
{\bf for $j:=2$ step $1$ until $n$ do}
\quad {\bf begin} ${\it accum}:=A[j]$; $k:=j-1$; $A[0]:=\it accum$;
\quad {\bf while $A[k]>\it accum$ do}
\qquad {\bf begin} $A[k+1]:=A[k]$; $k:=k-1$;
\qquad {\bf end};
\quad $A[k+1]:=\it accum$;
\quad {\bf end}.
}$$
\answer |{\obeylines \sfcode`;=3000|^^|\sfcode|\parbreak
|{\bf for $j:=2$ step $1$ until $n$ do}|\parbreak
|\quad {\bf begin} ${\it accum}:=A[j]$; $k:=j-1$; $A[0]:=\it accum$;|\parbreak
|\quad {\bf while $A[k]>\it accum$ do}|\parbreak
|\qquad {\bf begin} $A[k+1]:=A[k]$; $k:=k-1$;|\parbreak
|\qquad {\bf end};|\parbreak
|\quad $A[k+1]:=\it accum$;|\parbreak
|\quad {\bf end}.\par}|\par
\smallskip\noindent
[This is something like the ``poetry'' example in Chapter~14, but much
more difficult. Some manuals of style say that ^{punctuation} should inherit
the font of the preceding character, so that three kinds of semicolons
should be typeset; e.g., these experts recommend `$k:=j-1$; \
$A[0]:={}${\it accum;} \ {\bf end;}'. The author heartily disagrees.]

\subsection Spacing between formulas. ^{Displays} often contain more than one
formula; for example, an equation is frequently accompanied by a ^{side
condition}:
$$F_n=F_{n-1}+F_{n-2},\qquad n\ge2.$$
In such cases you need to tell \TeX\ how much space to put after the comma,
because \TeX's normal spacing conventions would bunch things together;
without special precautions you would get
$$F_n=F_{n-1}+F_{n-2}, n\ge2.$$

The traditional hot-metal technology for printing has led to some ingrained
standards for situations like this, based on what printers call a ``^{quad}''
of space. Since these standards seem to work well in practice, \TeX\ makes
it easy for you to continue the tradition: When you type `^|\quad|' in plain
\TeX\ format, you get a printer's quad of space in the horizontal direction.
Similarly, `^|\qquad|' gives you a double quad (twice as much); this
is the normal spacing for situations like
the $F_n$ example above. Thus, the recommended procedure is to type
\begintt
$$ F_n = F_{n-1} + F_{n-2}, \qquad  n \ge 2. $$
\endtt
It is perhaps worth reiterating that \TeX\ ignores all the spaces in math
mode (except, of course, the space after `|\qquad|', which is needed
to distinguish between `|\qquad|~|n|' and `|\qquadn|'); so the same result
would be obtained if you were to leave out all but one space:
\begintt
$$F_n=F_{n-1}+F_{n-2},\qquad n\ge2.$$
\endtt
Whenever you want spacing that differs from the normal conventions, you must
specify it explicitly by using control sequences such as |\quad| and |\qquad|.

\danger A quad used to be a square piece of blank type, $1\em$ wide and $1\em$
tall---approximately the size of a capital M, as explained in Chapter~10. This
tradition has not been fully retained: The control sequence |\quad| in plain
\TeX\ is simply an abbreviation for `|\hskip|~|1|^|em||\relax|', so \TeX's
quad has width but no height.

\danger You can use |\quad| in text as well as in formulas; for example,
Chapter~14 illustrates how |\quad| applies to poetry. When |\quad| appears
in a formula it stands for one~em in the current text font, independent of the
current math size or style or family. Thus, for example, |\quad| is just
as wide in a subscript as it is on the main line of a formula.

Sometimes a careless author will put two formulas next to each other in
the text of a paragraph. For example, you might find a sentence like this:
\begindisplay
The ^{Fibonacci} numbers satisfy $F_n=F_{n-1}+F_{n-2}$, \ $n\ge2$.
\enddisplay
Everybody who teaches proper ^{mathematical} ^{style} is agreed that formulas
ought to be separated by words, not just by commas; the author of that
sentence should at least have said `for $n\ge2$', not simply `$n\ge2$'.
But alas, such lapses are commonplace, and many prominent mathematicians
are hopelessly addicted to clusters of formulas. If we are not allowed to change
their writing style, we can at least insert extra space where they
neglected to insert an appropriate word. An additional interword space
generally works well in such cases; for example, the sentence above was
typeset thus:
\begintt
... $F_n=F_{n-1}+F_{n-2}$, \ $n\ge2$.}$$
\endtt
The `|\|\]' ^^{control space} here gives a visual separation that
partly compensates for the bad style.

\exercise Put the following paragraph into \TeX\ form, treating punctuation
and spacing carefully; also insert ^{ties} to prevent bad line breaks.
\begindisplay\baselineskip13pt
\vbox{\raggedright\hsize=310pt\parindent=0pt
Let $H$~be a Hilbert space, \
$C$~a closed bounded convex subset of~$H$, \
$T$~a nonexpansive self map of~$C$.
Suppose that as $n\to\infty$, \ $a_{n,k}\to0$ for each~$k$,
and $\gamma_n=\sum_{k=0}^\infty(a_{n,k+1}-a_{n,k})^+\to0$.
Then for each $x$~in~$C$, \
$A_nx=\sum_{k=0}^\infty a_{n,k}T^kx$ converges weakly
to a fixed point of~$T$.
} % taken from Bull. AMS 82 (1976), p 959; chosen by AMS in '78 for demo
\enddisplay
\answer |Let $H$~be a Hilbert space, \
$C$~a closed bounded convex subset of~$H$, \
$T$~a nonexpansive self map of~$C$.
Suppose that as $n\to\infty$, \ $a_{n,k}\to0$ for each~$k$,
and $\gamma_n=\sum_{k=0}^\infty(a_{n,k+1}-|\allowbreak|a_{n,k})^+\to0$.
 Then for each $x$~in~$C$,  \
$A_nx=\sum_{k=0}^\infty a_{n,k}T^kx$ converges weakly
to a fixed point of~$T$.|\par
[If any mathematicians are reading this, they might either appreciate
or resent the following attempt to edit the given paragraph
into a more acceptable style: ``%
Let $C$~be a closed, bounded, convex subset of a Hilbert space~$H$,
and let $T$~be a nonexpansive self map of~$C$.
Suppose that as $n\to\infty$, we have $a_{n,k}\to0$ for each~$k$,
and $\gamma_n=\sum_{k=0}^\infty(a_{n,k+1}-a_{n,k})^+\to0$.
Then for each $x$~in~$C$, the infinite sum
$A_nx=\sum_{k=0}^\infty a_{n,k}T^kx$ converges weakly
to a fixed point of~$T$.'']

\subsection Spacing within formulas. Chapter 16 says that \TeX\ does
automatic ^{spacing} of math formulas so that they look right, and this is
almost true. But occasionally you must give \TeX\ some help. The number of
possible math formulas is vast, and \TeX's spacing rules are rather
simple, so it is natural that exceptions should arise.  Of course, it is
desirable to have fine units of spacing for this purpose, instead of the
big chunks that arise from |\|\], |\quad|, and |\qquad|.

The basic elements of space that \TeX\
puts into formulas are called {\sl ^{thin spaces}}, {\sl ^{medium
spaces}}, and {\sl ^{thick spaces}}.  In order to get a feeling for these
units, let's take a look at the $F_n$ example again: Thick spaces occur
just before and after the = sign, and also before and after the $\ge$\thinspace;
medium spaces occur just before and after the $+$ sign.  Thin spaces are
slightly smaller, but noticeable; it's a thin space that makes the
difference between `loglog' and `$\log\log$'. The normal space between
words of a paragraph is approximately equal to two thin spaces.

\TeX\ inserts thin spaces, medium spaces, and thick spaces into formulas
automatically, but you can add your own spacing whenever you want to,
by using the control sequences ^^|\,|^^|\!|^^|\;|^^|\>|
$$\halign{\indent#\hfil&\quad#\hfil\cr
|\,|&thin space \ (normally 1/6 of a quad);\cr
|\>|&medium space \ (normally 2/9 of a quad);\cr
|\;|&thick space \ (normally 5/18 of a quad);\cr
|\!|&negative thin space \ (normally $-1/6$ of a quad).\cr}$$
In most cases you can rely on \TeX's spacing while you are typing a manuscript,
and you'll want to insert or delete space with these four control sequences
only in rare circumstances after you see what comes out.

\ddanger We observed a minute ago that |\quad| spacing does not
change with the style of formula, nor does it depend on the math font
families that are being used. But thin spaces, medium spaces, and thick
spaces do get bigger and smaller as the size of type gets bigger and
smaller; this is because they are defined in terms of ^\<muglue>, a~special
brand of glue intended for math spacing. You specify \<muglue> just
as if it were ordinary glue, except that the units are given in terms of
`^|mu|' (math units) instead of~|pt| or~|cm| or something else. For
example, Appendix~B contains the definitions
\begintt
\thinmuskip = 3mu
\medmuskip = 4mu plus 2mu minus 4mu
\thickmuskip = 5mu plus 5mu
\endtt
^^|\thinmuskip|^^|\medmuskip|^^|\thickmuskip|
and this defines the thin, medium, and thick spaces that \TeX\ inserts
into formulas. According to these specifications, thin spaces in plain
\TeX\ do not stretch or shrink; medium spaces can stretch a little, and
they can shrink to zero; thick spaces can stretch a lot, but they never shrink.

\ddanger There are 18 mu to an em, where the em is taken from family~2
(the math symbols family). In other words, ^|\textfont|~|2| defines the em
value for |mu| in display and text styles; ^|\scriptfont|~|2| defines the
em for script size material; and ^|\scriptscriptfont|~|2| defines it for
scriptscript size.

\ddanger You can insert math glue into any formula just by giving
the command `^|\mskip|\<muglue>'. For example, `|\mskip 9mu plus 2mu|'
inserts one half em of space, in the current size, together with some
stretchability. Appendix~B defines `|\,|' to be an abbreviation for
`|\mskip\thinmuskip|'.  Similarly, you can use the command `^|\mkern|'
when there is no stretching or shrinking; `|\mkern18mu|' gives one em of
horizontal space in the current size. \TeX\ insists that |\mskip| and
|\mkern| be used only with |mu|; conversely, ^|\hskip| and ^|\kern| (which
are also allowed in formulas) must never give units in |mu|.

Formulas involving ^{calculus} look best when an extra thin space appears
before $dx$ ^^{dx} or~$dy$ or~$d\,$whatever; but \TeX\ doesn't do this
automatically. Therefore a well-trained typist will remember to insert
`|\,|'  in examples like the following:
\beginmathdemo
\it Input&\it Output\cr
\noalign{\vskip2pt}
|$\int_0^\infty f(x)\,dx$|&\int_0^\infty f(x)\,dx\cr
|$y\,dx-x\,dy$|&y\,dx-x\,dy\cr
|$dx\,dy=r\,dr\,d\theta$|&dx\,dy=r\,dr\,d\theta\cr
|$x\,dy/dx$|&x\,dy/dx\cr \endmathdemo Notice that no `|\,|' was desirable
after the `|/|' in the last example.  Similarly, there's no need for
`|\,|' in cases like
\begindisplaymathdemo
|$$\int_1^x{dt\over t}$$|&\int_1^x{dt\over t}\cr
\endmathdemo
since the $dt$ appears all by itself in the numerator of a fraction; this
detaches it visually from the rest of the formula.

\exercise Explain how to handle the display
$$\int_0^\infty{t-ib\over t^2+b^2}e^{iat}\,dt=e^{ab}E_1(ab),\qquad a,b>0.$$
\answer |$$\int_0^\infty{t-ib\over t^2+b^2}e^{iat}\,dt=|\parbreak
        |    e^{ab}E_1(ab),\qquad a,b>0.$$|

\danger When physical ^{units} appear in a formula, they should be set in roman
type and separated from the preceding material by a thin space:
\beginmathdemo
|$55\rm\,mi/hr$|&55\rm\,mi/hr\cr
|$g=9.8\rm\,m/sec^2$|&g=9.8\rm\,m/sec^2\cr
|$\rm1\,ml=1.000028\,cc$|&\rm1\,ml=1.000028\,cc\cr
\endmathdemo

\dangerexercise Typeset the following display, assuming that `^|\hbar|'
generates `$\hbar$':
$$\hbar=1.0545\times10^{-27}\rm\,erg\,sec.$$
\answer |$$\hbar=1.0545\times10^{-27}\rm\,erg\,sec.$$|

\danger Thin spaces should also be inserted after ^{exclamation points}
(which stand for the ``^{factorial}'' operation in a formula), if the next
character is a letter or a number or an opening delimiter:
\beginmathdemo
|$(2n)!/\bigl(n!\,(n+1)!\bigr)$|&(2n)!/\bigl(n!\,(n+1)!\bigr)\cr
\noalign{\vskip6pt}
|$${52!\over13!\,13!\,26!}$$|&\displaystyle{52!\over13!\,13!\,26!}\cr
\endmathdemo

Besides these cases, you will occasionally encounter formulas in which
the symbols are bunched up too tightly, or where too much white space
appears, because of certain unlucky combinations of shapes. It's usually
impossible to anticipate optical glitches like this until you see the first
proofs of what you have typed; then you get to use your judgment about how
to add finishing touches that provide extra beauty, clarity, and finesse.
A tastefully applied `|\,|' or `|\!|'\ will open things up or close things
together so that the reader won't be distracted from the mathematical
significance of the formula. ^{Square root} signs and ^{multiple integrals} are
often candidates for such fine tuning. Here are some examples of situations
to look out for:
\beginmathdemo
|$\sqrt2\,x$|&\sqrt2\,x\cr
|$\sqrt{\,\log x}$|&\sqrt{\,\log x}\cr
|$O\bigl(1/\sqrt n\,\bigr)$|&O\bigl(1/\sqrt n\,\bigr)\cr
|$[\,0,1)$|&[\,0,1)\cr
|$\log n\,(\log\log n)^2$|&\log n\,(\log\log n)^2\cr
|$x^2\!/2$|&x^2\!/2\cr
|$n/\!\log n$|&n/\!\log n\cr
|$\Gamma_{\!2}+\Delta^{\!2}$|&\Gamma_{\!2}+\Delta^{\!2}\cr
|$R_i{}^j{}_{\!kl}$|&R_i{}^j{}_{\!kl}\cr
|$\int_0^x\!\int_0^y dF(u,v)$|&\int_0^x\!\int_0^y dF(u,v)\cr
\noalign{\vskip6pt}
|$$\int\!\!\!\int_D dx\,dy$$|&\displaystyle{\int\!\!\!\int_D dx\,dy}\cr
\endmathdemo
^^|\Gamma|^^|\Delta|^^|\intint|
In each of these formulas the omission of\/ |\,| or |\!|\ would lead to
somewhat less satisfactory results.

\ddanger Most of these examples where thin-space corrections are desirable
arise because of chance coincidences. For example, the superscript in
|$x^2/2$| leaves a hole before the slash ($x^2/2$); a negative thin
space helps to fill that hole. The positive thin space in |$\sqrt{\,\log x}$|
compensates for the fact that `$\log x$' begins with a tall, unslanted
letter; and so on. But two of the examples involve corrections that
were necessary because \TeX\ doesn't really know a great deal about
mathematics: \ (1)~In the formula
|$\log n(\log\log n)^2$|, \TeX\ inserts no thin space before the left
parenthesis, because there are similar formulas like |$\log n(x)$| where
no such space is desired. \ (2)~In the formula |$n/\log n$|, \TeX\
automatically inserts an unwanted thin~space before |\log|, since the slash is
treated as an ordinary symbol, and since a~thin space is usually desirable
between an ordinary symbol and an operator like |\log|.

\ddanger In fact, \TeX's rules for spacing in formulas are fairly simple.
A formula is converted to a math list as described at the end of Chapter~17,
and the math list consists chiefly of ``^{atoms}'' of eight basic types:
^{Ord}~(^{ordinary}), ^{Op}~(^{large operator}), ^{Bin}~(^{binary operation}),
^{Rel}~(^{relation}), ^{Open}~(^{opening}), ^{Close}~(^{closing}),
^{Punct}~(punctuation), ^^{punctuation} and ^{Inner}~(a delimited
subformula).  Other kinds of atoms, which arise from commands like
^|\overline| or ^|\mathaccent| or ^|\vcenter|, etc., are all treated as
type~Ord; subformulas delimited by ^|\left| and ^|\right|
are treated as type~Inner. The following table is
used to determine the spacing between pairs of adjacent atoms:
$$\baselineskip0pt\lineskip0pt
\halign to\hsize
 {\strut\hbox to\parindent{\it#\hfil}& % for the legend "Left atom"
  #\hfil\quad& % for the row labels
  #\hfil\tabskip 0pt plus 10pt& % for the rule at the left
  \hbox to 25pt{\tt\hss#\hss}& % for column 1
  \hbox to 25pt{\tt\hss#\hss}& % for column 2
  \hbox to 25pt{\tt\hss#\hss}& % for column 3
  \hbox to 25pt{\tt\hss#\hss}& % for column 4
  \hbox to 25pt{\tt\hss#\hss}& % for column 5
  \hbox to 25pt{\tt\hss#\hss}& % for column 6
  \hbox to 25pt{\tt\hss#\hss}& % for column 7
  \hbox to 25pt{\tt\hss#\hss}& % for column 8
  #\hfil\tabskip0pt\cr % for the rule at the right
&&&&\multispan7\hss\it Right atom\hss\cr
\noalign{\vskip3pt}
&&&\rm Ord&\rm Op&\rm Bin&\rm Rel&\rm Open&\rm Close&\rm Punct&\rm Inner\cr
\noalign{\vskip2pt}
\omit&&\multispan{10}\leaders\hrule\hfil\cr
\omit\vbox to 2pt{}&&\vrule&&&&&&&&&\vrule\cr
&Ord&\vrule&0&1&(2)&(3)&0&0&0&(1)&\vrule\cr
&Op&\vrule&1&1&*&(3)&0&0&0&(1)&\vrule\cr
&Bin&\vrule&(2)&(2)&*&*&(2)&*&*&(2)&\vrule\cr
Left&Rel&\vrule&(3)&(3)&*&0&(3)&0&0&(3)&\vrule\cr
atom&Open&\vrule&0&0&*&0&0&0&0&0&\vrule\cr
&Close&\vrule&0&1&(2)&(3)&0&0&0&(1)&\vrule\cr
&Punct&\vrule&(1)&(1)&*&(1)&(1)&(1)&(1)&(1)&\vrule\cr
&Inner&\vrule&(1)&1&(2)&(3)&(1)&0&(1)&(1)&\vrule\cr
\omit\vbox to 2pt{}&&\vrule&&&&&&&&&\vrule\cr
\omit&&\multispan{10}\leaders\hrule\hfil\cr}$$
^^{spacing table} ^^{math spacing table}
Here 0, 1, 2, and 3 stand for no space, thin space, medium space, and
thick space, respectively; the table entry is parenthesized if the space
is to be inserted only in display and text styles, not in script and
scriptscript styles. For example, many of the entries in the Rel row
and the Rel column are `{\tt(3)}'; this means that thick spaces are normally
inserted before and after relational symbols like `=', but not in
subscripts. Some of the entries in the table are `{\tt*}'; such cases
never arise, because Bin atoms must be preceded and followed by atoms
compatible with the nature of binary operations. Appendix~G contains
precise details about how math lists are converted to horizontal lists;
this conversion is done whenever \TeX\ is about to leave math mode, and the
interatomic spacing is inserted at that time.

\ddanger For example, the displayed formula specification
\begintt
$$x+y=\max\{x,y\}+\min\{x,y\}$$
\endtt
will be transformed into the sequence of atoms
\def\\#1{\vbox to 33pt{\vbox to 22pt{\vfill\hrule
      \hbox{\vrule\hskip-.4pt$#1$\hskip-.4pt\vrule}}\hrule\vfill}}%
\begindisplay
\vbox{\vskip-11pt\hbox{$
\\x\;\;\\+\;\;\\y\;\;\\=\;\;\\\max\;\;\\\{\;\;
\\x\;\;\\,\;\;\\y\;\;\\\}\;\;\\+\;\;\\\min\;\;\\\{
\;\;\\x\;\;\\,\;\;\\y\;\;\\\}$}\vskip-11pt}
\enddisplay
of respective types Ord, Bin, Ord, Rel, Op, Open, Ord, Punct, Ord, Close,
Bin, Op, Open, Ord, Punct, Ord, and Close.
Inserting spaces according to the table gives
$$\def\0{\thinspace}
\def\1{\thinspace{\tt\bslash,}\thinspace}
\def\2{\thinspace{\tt\bslash>}\thinspace}
\def\3{\thinspace{\tt\bslash;}\thinspace}
\halign{\indent\hfil#\cr
  Ord\2Bin\2Ord\3Rel\3Op\0Open\0Ord\0Punct\1Ord\0Close\2\qquad\cr
                 Bin\2Op\0Open\0Ord\0Punct\1Ord\0Close\cr}$$
and the resulting formula is
$$\vbox{\vskip-11pt\hbox{$
\\x\>\\+\>\\y\;\\=\;\\\max\\\{
\\x\\,\,\\y\\\}\>\\+\>\\\min\\\{
\\x\\,\,\\y\\\}$}\vskip-11pt}$$
i.e.,$$x+y=\max\{x,y\}+\min\{x,y\}\rlap{\quad.}$$
This example doesn't involve subscripts or superscripts; but subscripts and
superscripts merely get attached to atoms without changing the atomic type.

\ddangerexercise Use the table to determine what spacing \TeX\
will insert between the atoms of the formula `|$f(x,y)<x^2+y^2$|'.
\answer There are ten atoms (the first is $f$ and last is $y^2$); their types,
and the interatomic spacing, are respectively
\begindisplay \def\0{\thinspace}%
  \def\1{\thinspace{\tt\bslash,}\thinspace}%
  \def\2{\thinspace{\tt\bslash>}\thinspace}%
  \def\3{\thinspace{\tt\bslash;}\thinspace}
Ord\0Open\0Ord\0Punct\1Ord\0Close\3Rel\3Ord\2Bin\2Ord.
\enddisplay

\ddanger The plain \TeX\ macros ^|\bigl|, ^|\bigr|, ^|\bigm|, and ^|\big|
all produce identical delimiters; the only difference between them is that
they may lead to different spacing, because they make the delimiter into
different types of atoms: |\bigl| produces an Open atom, |\bigr| a~Close,
|\bigm| a~Rel, and |\big| an~Ord. On the other hand, when a subformula
appears between ^|\left| and ^|\right|, it is typeset by itself and placed
into an Inner atom. Therefore it is possible that a subformula enclosed by
|\left| and |\right| will be surrounded by more space than there would be
if that subformula were enclosed by |\bigl| and~|\bigr|. For example, Ord
followed by Inner (from |\left|) gets a thin space, but Ord followed by
Open (from |\bigl|) does not. The rules in Chapter~17 imply that the
construction `^|\mathinner||{\bigl({|\<subformula>|}\bigr)}|' within any
formula produces a result essentially equivalent to
`|\left(|\<subformula>|\right)|', when the \<subformula> doesn't end
with Punct, except that the ^{delimiters} are forced to
be of the |\big| size regardless of the height and depth of the subformula.

\danger \TeX's spacing rules sometimes fail when `\|' and `|\|\|'
appear in a formula, because $\vert$ and $\Vert$ are treated as ordinary symbols
^^{verticalline}^^|\verticalline|
instead of as delimiters. For example, consider the formulas
\beginlongmathdemo
|$|\||-x|\||=|\||+x|\||$|&\vert-x\vert=\vert+x\vert\cr
|$\left|\||-x\right|\||=\left|\||+x\right|\||$|&
  \left\vert-x\right\vert=\left\vert+x\right\vert\cr
|$\lfloor-x\rfloor=-\lceil+x\rceil$|&\lfloor-x\rfloor=-\lceil+x\rceil\cr
\endmathdemo
In the first case the spacing is wrong because \TeX\ thinks that the plus
sign is computing the sum of `$\vert$' and `$x$'. The use of\/ |\left| and
|\right| in the second example puts \TeX\ on the right track.  The third
example shows that no such corrections are needed with other delimiters,
because \TeX\ knows whether they are openings or closings.

\ddangerexercise Some perverse mathematicians use ^{brackets} backwards,
to denote ``^{open intervals}.''
Explain how to type the following bizarre formula: % from MR53 #3451
$\left]-\infty,T\right[\times\left]-\infty,T\right[$.
\answer |$\left]-\infty,T\right[\times\left]-\infty,T\right[$|. \ (Or one could
say ^|\mathopen| and ^|\mathclose| instead of\/ |\left| and |\right|;
then \TeX\ would not choose the size of the delimiters, nor would it consider
the subformulas to be of type Inner.) \
% that formula was quoted from MR review of paper by Mario Marino
% in Ricerche Mat. 24 (1975), no.~1, 152--171
Open intervals are more clearly expressed in print
by using parentheses instead of reversed brackets; for example,
compare `$(-\infty,T)\times(-\infty,T)$' to the given formula.

\ddangerexercise Study Appendix G and determine what spacing will be used
in the formula `|$x++1$|'. Which of the plus signs will be regarded
as a ^{binary operation}?
\answer The first |+| will become a Bin atom, the second an Ord; hence
the result is $x$, medium space, $+$, medium space, $+$, no space, 1.

\subsection Ellipses\/ {\rm(``three dots'')}. ^^{ellipses}
Mathematical copy looks much nicer if you are careful about how groups of
^{three dots} are typed in formulas and text. Although it looks fine to
type `|...|'\ on a typewriter that has fixed spacing, the result looks too
crowded when you're using a printer's fonts:  `|$x...y$|' results in
`$x...y$', and such close spacing is undesirable except in subscripts or
superscripts.

An ellipsis can be indicated by two different kinds of dots, one higher
than the other; the best mathematical traditions distinguish between
these two possibilities. It is generally correct to produce formulas like
\begindisplay
$\displaystyle x_1+\cdots+x_n\qquad {\rm and}\qquad (x_1,\ldots,x_n),$
\enddisplay
but wrong to produce formulas like
\begindisplay
$\displaystyle x_1+\ldots+x_n\qquad {\rm and}\qquad (x_1,\cdots,x_n).$
\enddisplay
The plain \TeX\ format of Appendix B allows you to solve the ``three dots''
problem very simply, and everyone will be envious of the beautiful formulas
that you produce. The idea is simply to type ^|\ldots| when you want three
low dots~($\,\ldots\,$), and ^|\cdots| when you want three vertically centered
dots~($\,\cdots\,$).

In general, it is best to use |\cdots| between $+$ and $-$ and~$\times$
signs, and also between $=$~signs or $\le$~signs or $\subset$ signs or other
similar relations. Low dots are used between ^{commas}, and when things
are juxtaposed with no signs between them at all. For example:
\beginmathdemo
|$x_1+\cdots+x_n$|&x_1+\cdots+x_n\cr
|$x_1=\cdots=x_n=0$|&x_1=\cdots=x_n=0\cr
|$A_1\times\cdots\times A_n$|&A_1\times\cdots\times A_n\cr
|$f(x_1,\ldots,x_n)$|&f(x_1,\ldots,x_n)\cr
|$x_1x_2\ldots x_n$|&x_1x_2\ldots x_n\cr
|$(1-x)(1-x^2)\ldots(1-x^n)$|&(1-x)(1-x^2)\ldots(1-x^n)\cr
|$n(n-1)\ldots(1)$|&n(n-1)\ldots(1)\cr
\endmathdemo

\exercise Type the formulas `$x_1+x_1x_2+\cdots+x_1x_2\ldots x_n$' and
`$(x_1,\ldots,x_n)\cdot(y_1,\ldots,y_n)=x_1y_1+\cdots+x_ny_n$'.
\ [{\sl Hint:\/} A single raised dot is called `^|\cdot|'.]
\answer |$x_1+x_1x_2+\cdots+x_1x_2\ldots x_n$| \ and\hfil\break
|$(x_1,\ldots,x_n)\cdot(y_1,\ldots,y_n)=x_1y_1+\cdots+x_ny_n$|.

But there's an important special case in which |\ldots| and |\cdots|
don't give the correct spacing, namely when they appear at the very end
of a formula, or when they appear just before a closing delimiter like~`|)|'.
In such situations an extra ^{thin space} is needed. For example, consider
sentences like this:
\begindisplay
Prove that $(1-x)^{-1}=1+x+x^2+\cdots\,$.\cr
Clearly $a_i<b_i$ for $i=1$,~2, $\ldots\,$,~$n$.\cr
The coefficients $c_0$,~$c_1$, \dots,~$c_n$ are positive.\cr
\enddisplay
To get the first sentence, the author typed
\begintt
Prove that $(1-x)^{-1}=1+x+x^2+\cdots\,$.
\endtt
Without the `^|\,|' the period would have come too close to the |\cdots|.
Similarly, the second sentence was typed thus:
\begintt
Clearly $a_i<b_i$ for $i=1$,~2, $\ldots\,$,~$n$.
\endtt
Notice the use of ^{ties}, which prevent bad line breaks as explained in
Chapter~14. Such ellipses are extremely common in some forms of mathematical
writing, so plain \TeX\ allows you to say just `^|\dots|' as an abbreviation
for `|$\ldots\,$|' in the text of a paragraph. The third sentence
can therefore be typed
\begintt
The coefficients $c_0$,~$c_1$, \dots,~$c_n$ are positive.
\endtt

\exercise B. C. ^{Dull} tried to take a shortcut by typing the
second example this way:
\begintt
Clearly $a_i<b_i$ for~$i=1, 2, \ldots, n$.
\endtt
What's so bad about that?
\answer The commas belong to the sentence, not to the formula; his
decision to put them into math mode meant that \TeX\ didn't put large
enough spaces after them. Also, his formula `$i=1, 2, \ldots, n$' allows
no breaks between lines, except after the $=$, so he's risking
overfull box problems. But suppose the sentence had been more terse:
\begindisplay
Clearly $a_i<b_i$ \ ($i=1,2,\ldots,n$).
\enddisplay
Then his idea would be basically correct:
\begintt
Clearly $a_i<b_i$ \ ($i=1,2,\ldots,n$).
\endtt

\exercise How do you think the author typed the ^{footnote} in Chapter 4
of this book?
\answer $\ldots$ |never\footnote*{Well \dots, hardly ever.} have| $\ldots$

\subsection Line breaking.  When you have formulas in a paragraph, \TeX\
may have to break them between lines.  This is a necessary evil, something
like the hyphenation of words; we want to avoid it unless the alternative
is worse. ^^{line breaking in math} ^^{breaking formulas between lines}

A formula will be broken only after a relation symbol like $=$
or~$<$ or $\to$, or after a binary operation symbol like $+$ or $-$ or
$\times$, where the relation or binary operation is on the ``outer level'' of
the formula (i.e., not enclosed in |{...}| and not part of an `|\over|'
construction). For example, if you type
\begintt
$f(x,y) = x^2-y^2 = (x+y)(x-y)$
\endtt
in mid-paragraph, there's a chance that \TeX\ will break after either of the
|=|~signs (it prefers this) or after the~|-| or~|+| or~|-| (in an emergency).
But there won't be a break after the comma in any case---commas
after which breaks are desirable shouldn't appear between |$|'s.

If you don't want to permit breaking in this example except after the
|=|~signs, you could type
\begintt
$f(x,y) = {x^2-y^2} = {(x+y)(x-y)}$
\endtt
because these additional braces ``freeze'' the ^{subformulas}, putting them
into unbreakable boxes in which the glue has been set to its natural width.
But it isn't necessary to bother worrying about such things unless \TeX\
actually does break a formula badly, since the chances of this are
pretty slim.

\danger A ``^{discretionary multiplication sign}'' is allowed in formulas:
If you type `|$(x+y)\*(x-y)$|', \TeX\ will treat the ^|\*| something like
the way it treats \hbox{|\-|}; namely, a line break will be allowed at
that place, with the hyphenation penalty. However, instead of inserting a
hyphen, \TeX\ will insert a $\times$ sign in text size.

\danger If you do want to permit a break at some point in the outer level
of a formula, you can say ^|\allowbreak|. For example, if the formula
\begintt
$(x_1,\ldots,x_m,\allowbreak y_1,\ldots,y_n)$
\endtt
appears in the text of a paragraph, \TeX\ will allow it to be broken into the
two pieces `$(x_1,\ldots,x_m,$' and `$y_1,\ldots,y_n)$'.

\ddanger The penalty for breaking after a Rel atom is called ^|\relpenalty|,
and the penalty for breaking after a Bin atom is called ^|\binoppenalty|.
Plain \TeX\ sets |\relpenalty=500| and |\binoppenalty=700|. You can change
the penalty for breaking in any particular case by typing `^|\penalty|\<number>'
immediately after the atom in question; then the number you have specified
will be used instead of the ordinary penalty. For example, you can prohibit
breaking in the formula `$x=0$' by typing `|$x=\nobreak0$|', since
^|\nobreak| is an abbreviation for `|\penalty10000 |'.

\ddangerexercise Is there any difference between the results of
`|$x=\nobreak0$|' and `|${x=0}$|'?
\answer Neither formula will be broken between lines, but the thick spaces
in the second formula will be set to their natural width while the thick
spaces in the first formula will retain their stretchability.

\ddangerexercise How could you prohibit all breaks in formulas, by making only
a few changes to the macros of plain \TeX?
\answer Set ^|\relpenalty||=10000| and ^|\binoppenalty||=10000|.
And you also need to change the definitions of\/ ^|\bmod| and ^|\pmod|,
which insert their own penalties.

\subsection Braces. A variety of different notations have sprung up involving
the symbols `$\{$' and `$\}$'; plain \TeX\ includes several control
sequences that help you cope with formulas involving such things.
^^{braces} ^^{leftbrace} ^^{rightbrace}

In simple situations, braces are used to indicate a ^{set} of objects;
for example, `$\{a,b,c\}$' stands for the set of three objects $a$, $b$,
and~$c$. There's nothing special about typesetting such formulas, except
that you must remember to use |\{| and |\}| for the braces:
^^|\leftbrace| ^^|\rightbrace|
\beginmathdemo
|$\{a,b,c\}$|&\{a,b,c\}\cr
|$\{1,2,\ldots,n\}$|&\{1,2,\ldots,n\}\cr
|$\{\rm red,white,blue\}$|&\{\rm red,white,blue\}\cr
\endmathdemo
A slightly more complex case arises when a set is indicated by giving a
generic element followed by a specific condition; for example, `$\{\,x\mid
x>5\,\}$' stands for the set of all objects $x$ that are greater than~5.
In such situations the control sequence ^|\mid| should be used for the
^{vertical bar}, and thin spaces should be inserted inside the braces:
\beginmathdemo
|$\{\,x\mid x>5\,\}$|&\{\,x\mid x>5\,\}\cr
|$\{\,x:x>5\,\}$|&\{\,x:x>5\,\}\cr
\endmathdemo
(Some authors prefer to use a ^{colon} instead of `$\mid$', as in the second
example here.) \ When the delimiters get larger, as in
\begindisplay
$\displaystyle\bigl\{\,\bigl(x,f(x)\bigr)\bigm\vert x\in D\,\bigr\}$
\enddisplay
they should be called ^|\bigl|, ^|\bigm|, and~^|\bigr|; for example,
the formula just given would be typed
\begintt
\bigl\{\,\bigl(x,f(x)\bigr)\bigm|char`||x\in D\,\bigr\}
\endtt
and formulas that involve still larger delimiters would use ^|\Big| or
^|\bigg| or~even ^|\Bigg|, as explained in Chapter~17.

\exercise How would you typeset the formula
$\bigl\{\,x^3\bigm\vert h(x)\in\{-1,0,+1\}\,\bigr\}$?
\answer |$\bigl\{\,x^3\bigm|\||h(x)\in\{-1,0,+1\}\,\bigr\}$|.

\dangerexercise Sometimes the condition that defines a set is given as
a fairly long English description, not as a formula; for example, consider
`$\{\,p\mid p$~and $p+2$ are prime$\,\}$'. An hbox would do the job:
\begintt
$\{\,p\mid\hbox{$p$ and $p+2$ are prime}\,\}$
\endtt
but a long formula like this is troublesome in a paragraph, since an hbox
cannot be broken between lines, and since the glue inside the |\hbox| does
not vary with the interword glue in the line that contains it. Explain how
the given formula could be typeset with line breaks allowed. [{\sl Hint:\/}
Go back and forth between math ^{mode} and horizontal mode.]
\answer |$\{\,p\mid p$~and $p+2$ are prime$\,\}$|, assuming that
^|\mathsurround| is zero. The more difficult alternative
`|$\{\,p\mid p\ {\rm and}\ p+2\rm\ are\ prime\,\}$|' is not a solution,
because line breaks do not occur at |\|\] ^^|\space|
(or at glue of any kind) within math formulas. Of course it may be best to
display a formula like this, instead of breaking it between lines.

Displayed formulas often involve another sort of brace, to indicate a choice
between various alternatives, as in the construction
\begindisplay
$\displaystyle\vert x\vert=\cases{x,&if $x\ge0$;\cr -x,&otherwise.\cr}$
\enddisplay
^^{selection, see cases} ^^{alternatives, see cases} ^^{choices, see cases}
You can typeset it with the control sequence ^|\cases|:
\begintt
$$|char`||x||=\cases{x,&if $x\ge0$;\cr
             -x,&otherwise.\cr}$$
\endtt
Look closely at this example and notice that it uses the character |&|,
^^{ampersand} which we said in Chapter~7 was reserved for special purposes.
Here for the first time in this manual we have an example of why |&|~is
so special: Each of the cases has two parts, and the~|&| separates those
parts. To the left of the~|&| is a math formula that is implicitly
enclosed in |$...$|; to the right of the~|&| is ordinary text, which is
{\sl not\/} implicitly enclosed in |$...$|. For example, the `|-x,|' in
the second line will be typeset in math mode, but the `|otherwise|' will
be typeset in horizontal mode. Blank spaces after the~|&| are ignored.
There can be any number of cases, but there usually are at least two.
Each case should be followed by ^|\cr|. Notice that the |\cases| construction
typesets its own `$\{$'; there is no corresponding `$\}$'.

\exercise Typeset the display \lower12pt\null\
$\smash{\displaystyle
f(x)=\cases{1/3&if $0\le x\le1$;\cr 2/3&if $3\le x\le4$;\cr 0&elsewhere.\cr}
}$
\answer |$$f(x)=\cases{1/3&if $0\le x\le1$;\cr 2/3&if $3\le x\le4$;\cr|\hfil
\break|0&elsewhere.\cr}$$|

\danger You can insert `^|\noalign||{|$\langle$vertical mode
material$\rangle$|}|' just after any \kern-1pt|\cr| within |\cases|, as
explained in Chapter~22, because |\cases| is an application of the general
alignment constructions considered in that chapter. For example, the
command `|\noalign{\vskip2pt}|' can be used to put a little extra space
between two of the cases.

\danger ^{Horizontal braces} will be set over or under parts of a displayed
formula if you use the control sequences ^|\overbrace| or ^|\underbrace|.
Such constructions are considered to be large operators like |\sum|, so you
can put limits above them or below them by specifying superscripts or
subscripts, as in the following examples:
\beginlongdisplaymathdemo
\noalign{\vskip9pt}
|$$\overbrace{x+\cdots+x}^{k\rm\;times}$$|&
  \overbrace{x+\cdots+x}^{k\rm\;times}\cr
\noalign{\vskip-6pt}
|$$\underbrace{x+y+z}_{>\,0}.$$|&
  \underbrace{x+y+z}_{>\,0}.\cr
\endmathdemo

\subsection Matrices. Now comes the fun part. Mathematicians in many different
disciplines like to construct rectangular arrays of formulas that have been
arranged in rows and columns; such an ^{array} is called a {\sl^{matrix}}.
Plain \TeX\ provides a ^|\matrix| control sequence that makes it convenient
to deal with the most common types of matrices.

For example, suppose that you want to specify the display
$$A=\left(\matrix{x-\lambda&1&0\cr
                  0&x-\lambda&1\cr
                  0&0&x-\lambda\cr}\right).$$
All you do is type
\begintt
$$A=\left(\matrix{x-\lambda&1&0\cr
                  0&x-\lambda&1\cr
                  0&0&x-\lambda\cr}\right).$$
\endtt
^^|\lambda|
This is very much like the |\cases| construction we looked at earlier;
each row of the matrix is followed by~|\cr|, and `|&|'~signs are used
between the individual entries of each row. Notice, however, that you are
supposed to put your own |\left| and |\right| delimiters around the matrix;
this makes |\matrix| different from |\cases|, which inserts a big `$\{$'
automatically.  The reason is that |\cases| always involves a left brace,
but different delimiters are used in different matrix constructions. On
the other hand, parentheses are used more often than other delimiters, so
you can write ^|\pmatrix| if you want plain \TeX\ to fill in the
parentheses for you; the example above then reduces to
\begintt
$$A=\pmatrix{x-\lambda&...&x-\lambda\cr}.$$
\endtt

\dangerexercise Typeset the display \ \lower12pt\null
$\tenpoint\smash{\displaystyle
\left\lgroup\matrix{a&b&c\cr d&e&f\cr}\right\rgroup
  \left\lgroup\matrix{u&x\cr v&y\cr w&z\cr}\right\rgroup
}$, \
using ^|\lgroup| and ^|\rgroup|.
\answer |$$\left\lgroup\matrix{a&b&c\cr d&e&f\cr}\right\rgroup|\hfil\break
|\left\lgroup\matrix{u&x\cr v&y\cr w&z\cr}\right\rgroup$$|.

\danger The individual entries of a matrix are normally centered in columns.
Each column is made as wide as necessary to accommodate the entries it
contains, and there's a ^{quad} of space between columns.
If you want to put something ^{flush right} in its column, precede it
by ^|\hfill|; if you want to put something ^{flush left} in its column,
follow it by~|\hfill|.

\danger Each entry of a matrix is treated separately from the others,
and it is typeset as a math formula in text style. Thus, for example,
if you say |\rm| in one entry, it does not affect the others.
Don't try to say `|{\rm x&y}|'.

Matrices often appear in the form of generic patterns that use ^{ellipses}
(i.e., dots) to indicate rows or columns that are left out. You can typeset
such matrices by putting the ellipses into rows and/or columns of their own.
Plain \TeX\ provides ^|\vdots| (vertical dots) and ^|\ddots| (diagonal dots)
as companions to ^|\ldots| for constructions like this. For example, the
^{generic matrix}
$$A=\pmatrix{a_{11}&a_{12}&\ldots&a_{1n}\cr
             a_{21}&a_{22}&\ldots&a_{2n}\cr
             \vdots&\vdots&\ddots&\vdots\cr
             a_{m1}&a_{m2}&\ldots&a_{mn}\cr}$$
is easily specified:
\begintt
$$A=\pmatrix{a_{11}&a_{12}&\ldots&a_{1n}\cr
             a_{21}&a_{22}&\ldots&a_{2n}\cr
             \vdots&\vdots&\ddots&\vdots\cr
             a_{m1}&a_{m2}&\ldots&a_{mn}\cr}$$
\endtt

\medskip
\exercise How can you get \TeX\ to produce the ^{column vector} ^^{vector}
\lower18pt\null\ $\smash{\displaystyle
\pmatrix{y_1\cr \vdots\cr y_k\cr}
}$\quad?
\answer |\pmatrix{y_1\cr \vdots\cr y_k\cr}|.

\danger Sometimes a matrix is bordered at the top and left by formulas
that give labels to the rows and columns. Plain \TeX\ provides a special
macro called ^|\bordermatrix| for this situation. For example, the display
$$\tenmath
M=\bordermatrix{&C&I&C'\cr C&1&0&0\cr I&b&1-b&0\cr C'&0&a&1-a\cr}$$
is obtained when you type
\begintt
$$M=\bordermatrix{&C&I&C'\cr
                C&1&0&0\cr  I&b&1-b&0\cr  C'&0&a&1-a\cr}$$
\endtt
The first row gives the upper labels, which appear above the big left
and right parentheses; the first column gives the left labels, which are
typeset flush left, just before the matrix itself. The first column in
the first row is normally blank. Notice that |\bordermatrix| inserts
its own parentheses, like |\pmatrix| does.

\danger It's usually inadvisable to put matrices into the text of a paragraph,
because they are so big that they are better displayed. But occasionally
you may want to specify a small matrix like $1\,1\choose0\,1$, which you can
^^|\choose| ^^{matrix, small}
typeset for example as `|$1\,1\choose0\,1$|'. Similarly, the small matrix
$\bigl({a\atop l}{b\atop m}{c\atop n}\bigr)$ can be typeset as
\begintt
$\bigl({a\atop l}{b\atop m}{c\atop n}\bigr)$
\endtt
^^|\atop| The |\matrix| macro does not produce small arrays of this sort.

\subsection Vertical spacing. If you want to tidy up an unusual formula,
you know already how to move things farther apart or closer together, by
using positive or negative thin spaces. But such spaces affect only the
horizontal dimension; what if you want something to be moved higher
or lower? That's an advanced topic.

\danger Appendix B provides a few macros that can be used to fool \TeX\
into thinking that certain formulas are larger or smaller than they really
are; such tricks can be used to move other parts of the formula up or down
or left or right. For example, we have already discussed the use of
^|\mathstrut| in Chapter~16 and ^|\strut| in Chapter~17; these invisible
boxes caused \TeX\ to put square root signs and the denominators of
continued fractions into different positions than usual.

\danger If you say `^|\phantom||{|\<subformula>|}|' in any formula, plain
\TeX\ will do its spacing as if you had said simply
`|{|\<subformula>|}|', but the subformula itself will be invisible. Thus,
for example, `|\phantom{0}2|' takes up just as much space as `|02|' in the
current style, but only the~|2| will actually appear on the page. If you
want to leave blank space for a ^{new symbol} that has exactly the same
size as $\sum$, but if you are forced to put that symbol in by hand for
some reason, `|\mathop{\phantom\sum}|' will leave exactly the right amount
of blank space.  \ (The `^|\mathop|' here makes this phantom behave like
|\sum|, i.e., as a large operator.)

\danger Even more useful than |\phantom| is ^|\vphantom|, which makes
an invisible box whose height and depth are the same as those of
the corresponding |\phantom|, but
the width is zero. Thus, |\vphantom| makes a vertical ^{strut} that can
increase a formula's effective height or depth. Plain \TeX\ defines
|\mathstrut| to be an abbreviation for `|\vphantom(|'. There's also
^|\hphantom|, which has the width of a |\phantom|, but its height
and depth are zero.

\danger Plain \TeX\ also provides `^|\smash||{|\<subformula>|}|', a macro
that yields the same result as `|{|\<subformula>|}|' but makes the height and
depth zero. By using both |\smash| and |\vphantom| you can typeset any
subformula and give it any desired nonnegative height and depth. For example,
\begintt
\mathop{\smash\limsup\vphantom\liminf}
\endtt
produces a large operator that says `$\limsup$', but its height and depth
are those of\/ ^|\liminf| (i.e., the depth is zero). ^^|\limsup|

\def\undertext#1{$\underline{\hbox{#1}}$}
\ddangerexercise If you want to underline some text, you could use a macro like
\begintt
\def\undertext#1{$\underline{\hbox{#1}}$}
\endtt
to do the job. \undertext{But} \undertext{this} \undertext{doesn't}
\undertext{always} \undertext{work} \undertext{right}. Discuss better
alternatives. ^^{underlined text}
\answer |\def|\stretch|\undertext|\stretch|#1{$\underline|\stretch
|{\smash|\stretch|{\hbox|\stretch|{#1}}}$}| will underline the
\def\undertext#1{$\underline{\smash{\hbox{#1}}}$}%
words and cross \undertext{through} the descenders; or you could insert
|\vphantom{y}| before the |\hbox|, thereby lowering all of the underlines
to a position below all descenders. Neither of these gives exactly what is
wanted. \ (See also ^|\underbar| in Appendix~B\null.) \ Underlining is actually
not very common in fine typography, since font changes usually work just
as well or better, when you want to emphasize something. If you really want
underlined text, it's best to have a special font in which all
the letters are underlined.

\ddanger You can also use ^|\raise| and ^|\lower| to adjust the vertical
positions of boxes in formulas. For example, the formula
`|$2^{\raise1pt\hbox{$\scriptstyle n$}}$|' will have its superscript~$n$
one point higher than usual ($2^{\raise1pt\hbox{$\scriptstyle n$}}$ instead of
$2^n$). Note that it was necessary to say ^|\scriptstyle| in this example,
since the contents of an ^|\hbox| will normally be in text style even when
that hbox appears in a superscript, and since |\raise| can be used only in
connection with a box.  This method of positioning is not used extremely
often, but it is sometimes helpful if the ^|\root| macro doesn't put its
argument in a suitable place. For example,
\begindisplay
|\root\raise|\<dimen>|\hbox{$\scriptscriptstyle|\<argument>|$}\of...|
\enddisplay
will move the argument up by a given amount.

\ddanger Instead of changing the sizes of subformulas, or using |\raise|,
you can also control vertical spacing by changing the parameters
that \TeX\ uses when it is converting math lists to horizontal lists.
These parameters are described in Appendix~G\null; you need to be careful when
changing them, because such changes are ^{global} (i.e., not local to groups).
Here is an example of how such a change might be made: Suppose that you
are designing a format for ^{chemical typesetting}, and that you expect to be
setting a lot of formulas like `$\rm Fe_2^{+2}Cr_2O_4$'. You may not like the
fact that the subscript in~$\rm Fe_2^{+2}$ is lower than the subscript
in~$\rm Cr_2$; and you don't want to force users to type monstrosities like
\begintt
$\rm Fe_2^{+2}Cr_2^{\vphantom{+2}}O_4^{\vphantom{+2}}$
\endtt
just to get the formula
$\rm Fe_2^{+2}Cr_2^{\vphantom{+2}}O_4^{\vphantom{+2}}$
with all subscripts at the same level. Well, all you need to do is set
`^|\fontdimen||16\tensy=2.7pt|' and `|\fontdimen17\tensy=2.7pt|', assuming
that ^|\tensy| is your main symbol font (|\textfont2|); this lowers all
normal ^{subscripts} to a position $2.7\pt$ below the baseline, which is
enough to make room for a possible superscript that contains a plus sign.
Similarly, you can adjust the positioning of ^{superscripts} by changing
|\fontdimen14\tensy|. There are parameters for the position of the ^{axis line},
the positions of ^{numerator} and ^{denominator} in a generalized ^{fraction},
the spacing above and below ^{limits}, the default ^{rule thickness}, and so
on. Appendix~G gives precise details.

\subsection Special features for math hackers. \TeX\ has a few more primitive
operations for math mode that haven't been mentioned yet. They are
occasionally useful if you are designing special formats.

\ddanger If a glue or kern specification is immediately preceded by
`^|\nonscript|', \TeX\ will not use that glue or kern in
script or scriptscript styles. Thus, for example, the sequence
`|\nonscript\;|' produces exactly the amount of space specified by
`{\tt(3)}' in the spacing table for mathematics that appeared earlier
in this chapter.

\ddanger Whenever \TeX\ has scanned a |$| and is about to read a math formula
that appears in text, it will first read another list of tokens that has
been predefined by the command ^|\everymath||={|\<token list>|}|. \ (This is
analogous to |\everypar|, which was described in Chapter~14.) \ Similarly,
you can say ^|\everydisplay||={|\<token list>|}| to predefine a list of tokens
for \TeX\ to read just after it has scanned an opening |$$|, i.e., just
before reading a formula that is to be displayed. With |\everymath| and
|\everydisplay|, you can set up special conventions that you wish to apply
to all formulas.

\subsection Summary. We have discussed more different kinds of formulas in
this chapter than you will usually find in any one book of mathematics.
If you have faithfully done the exercises so far, you can face almost
any formula with confidence.

\newcount\chalcount \chalcount=0
\outer\def\challenge{\d@nger\chall}
\outer\def\cchallenge{\dd@nger\chall}
\def\chall{\global\advance\chalcount by1
  \dexercise \hbox{Challenge number \the\chalcount:\enskip}\ignorespaces}

\danger But here are a few more exercises, to help you review what you
have learned. Each of the following ``challenge formulas'' illustrates one
or more of the principles already discussed in this chapter. The author
confesses that he is trying to trip you~up on several of these.
Nevertheless, if you try each one before looking at the answer, and if
you're alert for traps, you should find that these formulas provide a good
way to consolidate and complete your knowledge.

\challenge Explain how to type the phrase `$n^{\rm th}$ root', where
`$n^{\rm th}$' is treated as a mathematical formula with a superscript in
roman type.
\answer |$n^{\rm th}$ root|. \ (Incidentally, it is also acceptable
to type `|$n$th|', getting `$n$th', in such situations; the fact that
the $n$ is in italics distinguishes it from the suffix. Typed manuscripts
generally render this with a hyphen, but `$n$-th' is frowned on nowadays
when an italic~$n$ is available.) ^^{nth}

\challenge $\qquad\tenmath{\bf S^{\rm-1}TS=dg}(\omega_1,\ldots,\omega_n)
=\bf\Lambda$.
\answer |${\bf S^{\rm-1}TS=dg}(\omega_1,|\stretch|\ldots,|\stretch|\omega_n)
=\bf\Lambda$|.
\ $\bigl($Did you notice the difference between ^|\omega| ($\omega$)
and~|w| ($w$)?$\bigr)$

\challenge $\qquad\tenmath\Pr(\,m=n\mid m+n=3\,)$.
\answer |$\Pr(\,m=n\mid m+n=3\,)$|. \ (Analogous to a set.) ^^|\Pr|

\challenge $\qquad\tenmath\sin18^\circ={1\over4}(\sqrt5-1)$.^^{degrees}
\answer |$\sin18^\circ={1\over4}(\sqrt5-1)$|. ^^|\circ|

\challenge $\qquad\tenmath k=1.38065\times10^{-16}\rm\,erg\,K^{-1}$.
\answer |$k=1.38065\times10^{-16}\rm\,erg\,K^{-1}$|.

\challenge $\qquad\tenmath \bar\Phi\subset NL_1^*/N=\bar L_1^*
  \subseteq\cdots\subseteq NL_n^*/N=\bar L_n^*$.
\answer |$\bar\Phi\subset NL_1^*/N=\bar L_1^*|\parbreak
        |  \subseteq\cdots\subseteq NL_n^*/N=\bar L_n^*$|.

\challenge $\qquad\tenmath I(\lambda)=\int\!\!\int_Dg(x,y)e^{i\lambda h(x,y)}
  \,dx\,dy$. % cf. Math. Comp. 37 (1981), 509
\answer |$I(\lambda)=\int\!\!\int_Dg(x,y)e^{i\lambda h(x,y)}\,dx\,dy$|.\hfil
\break
(Although three |\!|'s work out best between consecutive integral signs in
displays, the text style seems to want only two.) ^^{double integral}
^^{integral, multiple}

\challenge $\qquad\tenmath
  \int_0^1\!\cdots\int_0^1f(x_1,\ldots,x_n)\,dx_1\ldots\,dx_n$.
\answer |$\int_0^1\!\cdots\int_0^1f(x_1,\ldots,x_n)\,dx_1\ldots\,dx_n$|.

\challenge Here's a display.
$$\tenmath x_{2m}\equiv\cases{Q(X_m^2-P_2W_m^2)-2S^2&($m$ odd)\cr
  \noalign{\vskip2pt}
  P_2^2(X_m^2-P_2W_m^2)-2S^2&($m$ even)\cr}\pmod N.$$
\answer |$$x_{2m}\equiv\cases{Q(X_m^2-P_2W_m^2)-2S^2&($m$ odd)\cr|\parbreak
        |       \noalign{\vskip2pt} % spread the lines apart a little|\parbreak
        |       P_2^2(X_m^2-P_2W_m^2)-2S^2&($m$ even)\cr}\pmod N.$$|

\challenge And another. % ACP Eq. 1.2.9--33
$$\tenmath (1+x_1z+x_1^2z^2+\cdots\,)\ldots(1+x_nz+x_n^2z^2+\cdots\,)
  ={1\over(1-x_1z)\ldots(1-x_nz)}.$$
\answer |$$(1+x_1z+x_1^2z^2+\cdots\,)\ldots(1+x_nz+x_n^2z^2+\cdots\,)|\parbreak
        |  ={1\over(1-x_1z)\ldots(1-x_nz)}.$$| \ (Notice the uses of\/ |\,|.)

\challenge And another. % Eq. 1.2.9--9
$$\tenmath \prod_{j\ge0}\biggl(\sum_{k\ge0}a_{jk}z^k\biggr)
  =\sum_{n\ge0}z^n\,\Biggl(\sum_
     {\scriptstyle k_0,k_1,\ldots\ge0\atop
      \scriptstyle k_0+k_1+\cdots=n}
   a_{0k_0}a_{1k_1}\ldots\,\Biggr).$$
\answer |$$\prod_{j\ge0}\biggl(\sum_{k\ge0}a_{jk}z^k\biggr)|\parbreak
        |  =\sum_{n\ge0}z^n\,\Biggl(\sum_|\parbreak
        |     {\scriptstyle k_0,k_1,\ldots\ge0\atop|\parbreak
        |      \scriptstyle k_0+k_1+\cdots=n}|\parbreak
        |   a_{0k_0}a_{1k_1}\ldots\,\Biggr).$$|\par
\nobreak\smallskip\noindent Some people would prefer to have the latter
parentheses larger; but |\left| and |\right| come out a bit too large in this
case. It's not difficult to define ^|\bigggl| and ^|\bigggr| macros, analogous
to the definitions of\/ |\biggl| and |\biggr| in Appendix~B.

\challenge And, % cf ACP vol1 p64
$$\tenmath {(n_1+n_2+\cdots+n_m)!\over n_1!\,n_2!\ldots n_m!}
  ={n_1+n_2\choose n_2}{n_1+n_2+n_3\choose n_3}
    \ldots{n_1+n_2+\cdots+n_m\choose n_m}.$$
\answer |$${(n_1+n_2+\cdots+n_m)!\over n_1!\,n_2!\ldots n_m!}|\parbreak
        |  ={n_1+n_2\choose n_2}{n_1+n_2+n_3\choose n_3}|\parbreak
        |    \ldots{n_1+n_2+\cdots+n_m\choose n_m}.$$|

\challenge Yet another display. % found in Chaundy et al
$$\tenmath \def\\#1#2{(1-q^{#1_#2+n})} % to save typing
\Pi_R{a_1,a_2,\ldots,a_M\atopwithdelims[]b_1,b_2,\ldots,b_N}
  =\prod_{n=0}^R{\\a1\\a2\ldots\\aM\over\\b1\\b2\ldots\\bN}.$$
\answer |$$\def\\#1#2{(1-q^{#1_#2+n})} % to save typing|\parbreak
        |\Pi_R{a_1,a_2,\ldots,a_M\atopwithdelims[]b_1,b_2,\ldots,b_N}|\parbreak
        |  =\prod_{n=0}^R{\\a1\\a2\ldots\\aM\over\\b1\\b2\ldots\\bN}.$$|
^^|\atopwithdelims|

\challenge And another.
$$\tenmath \sum_{p\rm\;prime}f(p)=\int_{t>1}f(t)\,d\pi(t).$$
\answer |$$\sum_{p\rm\;prime}f(p)=\int_{t>1}f(t)\,d\pi(t).$$|

\challenge Still another.
$$\tenmath \{\underbrace{\overbrace{\mathstrut a,\ldots,a}
      ^{k\;a\mathchar`'\rm s},
    \overbrace{\mathstrut b,\ldots,b}
      ^{l\;b\mathchar`'\rm s}}_{k+l\rm\;elements}\}.$$
\answer |$$\{\underbrace{\overbrace{\mathstrut a,\ldots,a}|\parbreak
        |      ^{k\;a\mathchar`'\rm s},|\parbreak
        |    \overbrace{\mathstrut b,\ldots,b}|\parbreak
        |      ^{l\;b\mathchar`'\rm s}}_{k+l\rm\;elements}\}.$$|\par
\smallskip\noindent Notice how ^{apostrophes} (instead of primes) were obtained.

\challenge Put a ^|\smallskip| between the rows of matrices in the
compound matrix ^^{compound matrix}
$$\tenmath \pmatrix{\pmatrix{a&b\cr c&d\cr}&
             \pmatrix{e&f\cr g&h\cr}\cr
           \noalign{\smallskip}
           0&\pmatrix{i&j\cr k&l\cr}\cr}.$$
\answer |$$\pmatrix{\pmatrix{a&b\cr c&d\cr}&|\parbreak
        |             \pmatrix{e&f\cr g&h\cr}\cr|\parbreak
        |           \noalign{\smallskip}|\parbreak
        |           0&\pmatrix{i&j\cr k&l\cr}\cr}.$$|

\challenge Make the columns ^{flush left} here. % cf Polya/Szego VII.43.2
$$\tenmath \det\left\vert\,\matrix{
  c_0&c_1\hfill&c_2\hfill&\ldots&c_n\hfill\cr
  c_1&c_2\hfill&c_3\hfill&\ldots&c_{n+1}\hfill\cr
  c_2&c_3\hfill&c_4\hfill&\ldots&c_{n+2}\hfill\cr
  \,\vdots\hfill&\,\vdots\hfill&
       \,\vdots\hfill&&\,\vdots\hfill\cr
  c_n&c_{n+1}\hfill&c_{n+2}\hfill&\ldots&c_{2n}\hfill\cr
  }\right\vert>0.$$
\answer |$$\det\left|\||\,\matrix{|\parbreak
        |  c_0&c_1\hfill&c_2\hfill&\ldots&c_n\hfill\cr|\parbreak
        |  c_1&c_2\hfill&c_3\hfill&\ldots&c_{n+1}\hfill\cr|\parbreak
        |  c_2&c_3\hfill&c_4\hfill&\ldots&c_{n+2}\hfill\cr|\parbreak
        |  \,\vdots\hfill&\,\vdots\hfill&|\parbreak
        |       \,\vdots\hfill&&\,\vdots\hfill\cr|\parbreak
        |  c_n&c_{n+1}\hfill&c_{n+2}\hfill&\ldots&c_{2n}\hfill\cr|\parbreak
        |  }\right|\||>0.$$|

\cchallenge The main problem here is to prime the $\sum$.
^^|\sum prime| ^^{=def}
$$\tenmath \mathop{{\sum}'}_{x\in A}f(x)\mathrel{\mathop=^{\rm def}}
  \sum_{\scriptstyle x\in A\atop\scriptstyle x\ne0}f(x).$$
\answer |$$\mathop{{\sum}'}_{x\in A}f(x)\mathrel{\mathop=^{\rm def}}|\parbreak
        |  \sum_{\scriptstyle x\in A\atop\scriptstyle x\ne0}f(x).$$|\par
\smallskip\noindent
This works because |{\sum}| is type Ord (so its superscript is not set
above), but ^|\mathop||{{\sum}'}| is type Op (so its subscript is set below).
The limits are centered on $\sum'$, however, not on $\sum$. If you don't
like that, the remedy is more difficult; one solution is to use
|\sumprime_{x\in A}| where ^|\sumprime| is defined as follows:
\par\nobreak\medskip
|\def\sumprime_#1{\setbox0=\hbox{$\scriptstyle{#1}$}|\parbreak
|  \setbox2=\hbox{$\displaystyle{\sum}$}|\parbreak
|  \setbox4=\hbox{${}'\mathsurround=0pt$}|\parbreak
|  \dimen0=.5\wd0 \advance\dimen0 by-.5\wd2|\parbreak
|  \ifdim\dimen0>0pt|\parbreak
|    \ifdim\dimen0>\wd4 \kern\wd4 \else\kern\dimen0\fi\fi|\parbreak
|  \mathop{{\sum}'}_{\kern-\wd4 #1}}|

\cchallenge You may be ready now for this display.
$$\tenmath 2\uparrow\uparrow k\mathrel{\mathop=^{\rm def}}
  2^{2^{2^{\cdot^{\cdot^{\cdot^2}}}}}
    \vbox{\hbox{$\Big\}\scriptstyle k$}\kern0pt}.$$
\answer |$$2\uparrow\uparrow k\mathrel{\mathop=^{\rm def}}|\parbreak
        |  2^{2^{2^{\cdot^{\cdot^{\cdot^2}}}}}|\parbreak
        |    \vbox{\hbox{$\Big\}\scriptstyle k$}\kern0pt}.$$|\par

\cchallenge And finally, when you have polished off all the other examples,
here's the ultimate test. Explain how to obtain the ^{commutative diagram}
% from Invent. Math. 70 (1982), 34 % with "typo" corrected 99.12.03
$$\tenmath\def\normalbaselines{\baselineskip20pt\lineskip1pt\lineskiplimit0pt }
\def\mapright#1{\smash{
    \mathop{\longrightarrow}\limits^{#1}}}
\def\mapdown#1{\Big\downarrow
  \rlap{$\vcenter{\hbox{$\scriptstyle#1$}}$}}
\matrix{\noalign{\vskip6pt}&&&&&&0\cr
  &&&&&&\mapdown{}\cr
  0&\mapright{}&{\cal O}_C&\mapright\iota&
    \cal E&\mapright\rho&\cal L&\mapright{}&0\cr
  &&\Big\Vert&&\mapdown\phi&&\mapdown\psi\cr
  0&\mapright{}&{\cal O}_C&\mapright\pi&
    \pi_*{\cal O}_D&\mapright\delta&
    R^1f_*{\cal O}_V(-D)&\mapright{}&0\cr
  &&&&&&\mapdown{\theta_i\otimes\gamma^{-1}}\cr
  &&&&&&\hidewidth R^1f_*\bigl({\cal O}
    _V(-iM)\bigr)\otimes\gamma^{-1}\hidewidth\cr
  &&&&&&\mapdown{}\cr
  &&&&&&0\cr\noalign{\vskip6pt}}$$
using ^|\matrix|. \ (Many of the entries are blank.)
\answer If you have to do a lot of commutative diagrams, you will want to
define some macros like those in the first few lines of this solution.
The ^|\matrix| macro resets the baselines to ^|\normalbaselines|, because
other commands like |\openup| might have changed them, so
we redefine |\normalbaselines| in this solution. Some of the things
shown here haven't been explained yet, but Chapter~22 will reveal all.
\smallskip
|$$\def\normalbaselines{\baselineskip20pt|\parbreak
|  \lineskip3pt \lineskiplimit3pt }|\parbreak
|\def\mapright#1{\smash{|\parbreak
|    \mathop{\longrightarrow}\limits^{#1}}}|\parbreak
|\def\mapdown#1{\Big\downarrow|\parbreak
|  \rlap{$\vcenter{\hbox{$\scriptstyle#1$}}$}}|\parbreak
|\matrix{&&&&&&0\cr|\parbreak
|  &&&&&&\mapdown{}\cr|\parbreak
|  0&\mapright{}&{\cal O}_C&\mapright\iota&|\parbreak
|    \cal E&\mapright\rho&\cal L&\mapright{}&0\cr|\parbreak
|  &&\Big\Vert&&\mapdown\phi&&\mapdown\psi\cr|\parbreak
|  0&\mapright{}&{\cal O}_C&\mapright\pi&|\parbreak
|    \pi_*{\cal O}_D&\mapright\delta&|\parbreak
|    R^1f_*{\cal O}_V(-D)&\mapright{}&0\cr|\parbreak
|  &&&&&&\mapdown{\theta_i\otimes\gamma^{-1}}\cr|\parbreak
|  &&&&&&\hidewidth R^1f_*\bigl({\cal O}|\parbreak
|    _V(-iM)\bigr)\otimes\gamma^{-1}|^|\hidewidth||\cr|\parbreak
|  &&&&&&\mapdown{}\cr|\parbreak
|  &&&&&&0\cr}$$|

\subsection Words of advice. The number of different notations is
enormous and still growing, so you will probably continue to find
new challenges as you continue to type mathematical papers. It's a
good idea to keep a personal notebook in which you record all of
the nonobvious formulas that you have handled successfully,
showing both the final output and what you typed to get it.
Then you'll be able to refer back to those solutions when you
discover that you need to do something similar, a few months later.

If you're a mathematician who types your own papers, you have now learned
^^{author, typesetting by}
how to get enormously complex formulas into print, and you can do so
without going through an intermediary who may somehow distort their
meaning. But please, don't get too carried away by your newfound talent;
the fact that you are able to typeset your formulas with \TeX\ doesn't
necessarily mean that you have found the best notation for communicating
with the readers of your work.  Some notations will be unfortunate even
when they are beautifully formatted.

\endchapter

Mathematicians are like Frenchmen:\/
% Die Mathematiker sind eine Art Franzosen:
whenever you say something to them, they translate it into their own language,
% redet man zu ihnen, so u"bersetzen sie es in ihre Sprache,
and at once it is something entirely different.
% und dann ist es alsobald ganz etwas Anders.
\author ^{GOETHE}, {\sl Maxims and Reflexions\/} (1829)
% see Schriften der Goethe-Gesellschaft, vol 21, pp 266 and 389

\bigskip

The best notation is no notation;
whenever it is possible to avoid the use of a complicated alphabetic apparatus,
avoid it.
A good attitude to the preparation of written mathematical exposition
is to pretend that it is spoken.
Pretend that you are explaining the subject to a friend
on a long walk in the woods, with no paper available;
fall back on symbolism only when it is really necessary.
\author PAUL ^{HALMOS}, {\sl How to Write Mathematics\/} (1970)
  % in {\sl L'Enseignement Math\'ematique\/}
  % vol 16, 123--152; section 15; reprinted in AMS pub "How to Write Math"

\eject
